% =====================================================================
% 系统开发工具基础 · 实验报告（第 1--4 题）
%
% 编译方式（在 WSL 中）：
%   latexmk -xelatex -halt-on-error 实验报告.tex
%
% 依赖安装（首次需要）：
%   sudo apt install -y texlive-latex-base texlive-latex-recommended \
%                        texlive-latex-extra latexmk texlive-lang-chinese \
%                        texlive-xetex fonts-noto-cjk poppler-utils
%   （Ubuntu 26.04 无 fandol 字体，ctex 使用 fontset=ubuntu + Noto CJK）
%
% 推荐编辑器：VS Code + LaTeX Workshop 扩展（配合 WSL 扩展在 WSL 中编译）
% =====================================================================
\documentclass[UTF8,fontset=ubuntu]{ctexart}

\usepackage{amsmath}
\usepackage{booktabs}
\usepackage{geometry}
\usepackage{listings}
\usepackage[hidelinks]{hyperref}

\geometry{a4paper, margin=2.5cm}

\lstset{
  basicstyle=\ttfamily\small,
  breaklines=true,
  frame=single,
  framerule=0.4pt,
  tabsize=4,
  showstringspaces=false,
  columns=fullflexible,
  keepspaces=true,
}

\title{\textbf{系统开发工具基础 · 实验报告（第 1--4 题）}}
\author{姓名：李云龙 \quad 学号：25030021043}
\date{\today}

\begin{document}

\maketitle

\begin{center}
\begin{tabular}{@{}p{3cm}p{11.5cm}@{}}
\toprule
姓名 & 李云龙 \\
学号 & 25030021043 \\
实验日期 & \today \\
实验环境 & Windows 11 + WSL2（Ubuntu）；bash + GNU coreutils（find / cp / sort / uniq / awk / head 等）；git；TeX Live（latexmk、pdflatex）；poppler-utils（pdftotext） \\
\bottomrule
\end{tabular}
\end{center}

\begin{quote}
\textbf{说明：}以下各题“运行结果”均为在 WSL（Ubuntu 26.04 LTS）中实际运行对应脚本得到的真实输出（运行目录：\lstinline|/mnt/e/Dyer/系统开发工具基础|）。各题产物均保存在 E 盘对应目录（q01/、q02/、q03/、q04/）。
\end{quote}

\section{实验目的}
\begin{enumerate}
\item 掌握含空格文件名、隐藏文件的操作方法（引号、\lstinline|find -exec|、\lstinline|cp --parents|），以及目录/文件权限设置。
\item 掌握 Shell 管道与 \lstinline|awk|、\lstinline|sort|、\lstinline|uniq|、\lstinline|head| 等文本处理工具，完成日志统计。
\item 掌握 Git 分支与合并操作，能够主动制造、解决并解释一次合并冲突。
\item 掌握 LaTeX 文档编辑（公式、表格、交叉引用）与 \lstinline|latexmk| 命令行构建 PDF。
\end{enumerate}

\section{实验一：含空格文件名的批量整理（q01）}

\subsection{完整脚本（scripts/q01.sh）}
\begin{lstlisting}[language=bash]
#!/bin/bash
set -e
SID=25030021043            # 学号
mkdir -p q01/input/docs q01/input/tmp
cd q01

# 1. 创建文件
printf 'alpha\nbeta\n' > 'input/docs/notes one.txt'
echo 'hidden' > input/docs/.secret.txt
touch input/tmp/empty.txt
printf '2026-08-24 10:00:00 INFO start\n2026-08-24 10:01:00 INFO stop\n' > input/run.log

# 2. 绝对路径 + 长格式列表（含隐藏文件）
pwd
ls -la input

# 3. 复制所有 .txt 并保留相对目录结构
mkdir -p "work/$SID"
(cd input && find . -type f -name '*.txt' -exec cp --parents {} "../work/$SID/" \;)

# 4. 权限：目录 750，普通文件 640
find "work/$SID" -type d -exec chmod 750 {} +
find "work/$SID" -type f -exec chmod 640 {} +

# 5. inventory.txt：相对路径 + 字节数
find "work/$SID" -type f -exec wc -c {} \; > inventory.txt

# 验证
ls -laR "work/$SID"   # -a 显示隐藏文件；目录 drwxr-x---，文件 -rw-r-----
cat inventory.txt
\end{lstlisting}

\subsection{运行结果}
\textbf{pwd}（q01 的绝对路径）：
\begin{lstlisting}
/mnt/e/Dyer/系统开发工具基础/q01
\end{lstlisting}

\textbf{ls -la input}（长格式，含隐藏项目；ls -laR input 可连带列出子目录内容）：
\begin{lstlisting}
total 0
drwxr-xr-x 1 root root 4096 Aug 24 21:04 .
drwxr-xr-x 1 root root 4096 Aug 24 21:04 ..
drwxr-xr-x 1 root root 4096 Aug 24 21:04 docs
-rw-r--r-- 1 root root   61 Aug 24 21:04 run.log
drwxr-xr-x 1 root root 4096 Aug 24 21:04 tmp
\end{lstlisting}

\textbf{复制后的结构与权限}（ls -laR work/25030021043）：
\begin{lstlisting}
work/25030021043/:
total 0
drwxr-x--- 1 root root 4096 Aug 24 21:04 .
drwxr-xr-x 1 root root 4096 Aug 24 21:04 ..
drwxr-x--- 1 root root 4096 Aug 24 21:04 docs
drwxr-x--- 1 root root 4096 Aug 24 21:04 tmp

work/25030021043/docs:
total 0
drwxr-x--- 1 root root 4096 Aug 24 21:04 .
drwxr-x--- 1 root root 4096 Aug 24 21:04 ..
-rw-r----- 1 root root    7 Aug 24 21:04 .secret.txt
-rw-r----- 1 root root   11 Aug 24 21:04 notes one.txt

work/25030021043/tmp:
total 0
drwxr-x--- 1 root root 4096 Aug 24 21:04 .
drwxr-x--- 1 root root 4096 Aug 24 21:04 ..
-rw-r----- 1 root root    0 Aug 24 21:04 empty.txt
\end{lstlisting}

\textbf{inventory.txt}（相对路径 + 字节数）：
\begin{lstlisting}
7 work/25030021043/docs/.secret.txt
11 work/25030021043/docs/notes one.txt
0 work/25030021043/tmp/empty.txt
\end{lstlisting}

\subsection{要点分析}
\begin{itemize}
\item \textbf{空格处理}：\lstinline|find -exec| 会为 \lstinline|{}| 自动加引号、逐个传参，\lstinline|cp --parents| 原样保留 \lstinline|notes one.txt| 中的空格，不会把文件名拆成两个参数。
\item \textbf{复制范围}：\lstinline|.secret.txt| 以 .txt 结尾，属于“所有 .txt 文件”，被一并复制；\lstinline|run.log| 不是 .txt，不复制；\lstinline|tmp/empty.txt| 是空文件但后缀符合，也复制。
\item \textbf{目录结构}：\lstinline|find . -type f| 输出 \lstinline|./docs/...| 形式路径，\lstinline|cp --parents| 在目标下重建 docs/、tmp/ 子目录，\lstinline|./| 会被路径解析自动归一化，不影响结构。
\item \textbf{权限}：目录 750 = rwxr-x---（所有者全部权限，组可读可执行，其他无）；普通文件 640 = rw-r-----（所有者读写，组只读，其他无）。
\item \textbf{字节数核对}：\texttt{alpha\textbackslash nbeta\textbackslash n} 为 5+1+4+1=11 字节；\texttt{hidden\textbackslash n} 为 7 字节；空文件 0 字节。
\end{itemize}

\section{实验二：随机访问日志统计（q02）}

\subsection{access.csv}
按题目给定内容创建（1 行表头 + 8 行数据，共 9 行）。

\subsection{analyze.sh（scripts/q02.sh 内生成）}
\begin{lstlisting}[language=bash]
#!/bin/bash
# 用法: ./analyze.sh <csv文件>
if [ $# -ne 1 ]; then
    echo "用法: $0 <csv文件>" >&2
    exit 2
fi
csv="$1"

if [ ! -f "$csv" ]; then
    echo "错误: 文件 $csv 不存在" >&2
    exit 1
fi

# 1) 5xx 数量最多的前 2 个 path（次数降序，相同按字典序）
tail -n +2 "$csv" | awk -F, '$4 >= 500 && $4 < 600 {print $3}' \
    | sort | uniq -c | sort -k1,1nr -k2,2 | head -n 2

# 2) 全部数据行的平均 latency_ms（表头不计入），保留两位小数
tail -n +2 "$csv" | awk -F, '{sum += $5; n++} END {printf "平均 latency: %.2f\n", sum/n}'
\end{lstlisting}

\subsection{运行结果}
对 access.csv 运行：
\begin{lstlisting}
      3 /api/orders
      2 /api/users
平均 latency: 193.75
\end{lstlisting}
（\lstinline|uniq -c| 输出的次数带右对齐前导空格，属正常现象）

对不存在的文件运行（错误进入 stderr，退出码非零）：
\begin{lstlisting}
错误: 文件 missing.csv 不存在      （stderr）
exit=1
\end{lstlisting}

\subsection{要点分析}
\begin{itemize}
\item \textbf{管道分解}：\lstinline|tail -n +2| 跳过表头 -> \lstinline|awk| 过滤状态码 500--599 并取出第 3 列 path -> \lstinline|sort| 让相同 path 相邻 -> \lstinline|uniq -c| 计数 -> \lstinline|sort -k1,1nr -k2,2| 按次数降序、次数相同按 path 字典序 -> \lstinline|head -n 2| 取前 2。
\item \textbf{5xx 计数核对}：/api/orders 3 次（500、502、500）、/api/users 2 次（503、500）、/login 1 次（500），故前 2 名为 /api/orders、/api/users。
\item \textbf{平均延迟核对}：$(120+310+180+90+260+20+420+150)/8 = 1550/8 = 193.75$，\lstinline|printf "%.2f"| 保留两位小数；表头被 \lstinline|tail -n +2| 排除，不计入。
\item \textbf{错误处理}：文件不存在时错误信息用 \lstinline|>&2| 写入 stderr，\lstinline|exit 1| 返回非零退出码，满足题目要求。
\end{itemize}

\section{实验三：制造、解决并解释一次合并冲突（q03）}

\subsection{操作步骤（scripts/q03.sh）}
\begin{enumerate}
\item \lstinline|git init|，将默认分支命名为 main，创建 config.txt（mode=normal）并完成初始提交。
\item 新建分支 feature-a，把内容改为 mode=safe 并提交。
\item 回到 main（此时仍在初始提交），从初始 main 新建分支 feature-b，把内容改为 mode=fast 并提交。
\item 回到 main，先 \lstinline|git merge feature-a|（快进合并，无冲突）。
\item 再 \lstinline|git merge feature-b|，产生冲突。
\end{enumerate}

\subsection{冲突原因}
feature-a 与 feature-b 都从初始提交（mode=normal）出发，各自把同一文件\textbf{同一行}改成了不同内容（mode=safe / mode=fast）。合并时 Git 找不到共同基线下的唯一答案，无法自动合并，于是产生冲突。冲突现场 config.txt 内容：
\begin{lstlisting}
<<<<<<< HEAD
mode=safe
=======
mode=fast
>>>>>>> feature-b
\end{lstlisting}

\subsection{冲突解决}
手工编辑（保留 mode=safe，另加一行 note=reviewed），然后标记已解决并提交，完成本次合并：
\begin{lstlisting}[language=bash]
printf 'mode=safe\nnote=reviewed\n' > config.txt
git add config.txt
git commit -m "merge feature-b: keep mode=safe, add note=reviewed"
\end{lstlisting}

\subsection{验证结果}
\textbf{git status}：\lstinline|nothing to commit, working tree clean|（工作区干净）。

\textbf{最终 config.txt}：
\begin{lstlisting}
mode=safe
note=reviewed
\end{lstlisting}

\textbf{git log --all --graph --decorate --oneline}：
\begin{lstlisting}
*   8703ad7 (HEAD -> main) merge feature-b: keep mode=safe, add note=reviewed
|\
| * cac5500 (feature-b) feature-b: set mode=fast
* | e3b1f51 (feature-a) feature-a: set mode=safe
|/
* 733bff0 init: mode=normal
\end{lstlisting}

\subsection{小结}
\begin{itemize}
\item 合并冲突的本质：双方从同一基线修改了同一区域，Git 无法判断保留哪一份，需要人工决策。
\item 解决流程：编辑冲突文件（删除 \lstinline|<<<<<<<|、\lstinline|=======|、\lstinline|>>>>>>>| 标记）-> \lstinline|git add| -> \lstinline|git commit|；也可用 \lstinline|git checkout --ours/--theirs| 直接选取一侧。
\end{itemize}

\section{实验四：修复并构建一页技术说明（q04）}

\subsection{任务与完整代码}
按题目模板补全内容：加入一个带编号和 label 的公式并在正文中引用（\lstinline|\eqref|）；加入一个 2 行 2 列表格，设置 caption 与 label 并在正文中引用（\lstinline|\ref|）；最后用 \lstinline|latexmk -pdf -halt-on-error| 构建 PDF，确认交叉引用不显示 ``??''。report.tex 完整内容如下：
\begin{lstlisting}
\documentclass{article}
\usepackage{amsmath}
\begin{document}
\title{Tool Report}
\author{Li Yunlong -- 25030021043}
\maketitle
\section{Result}

The mass--energy equivalence is given by Equation~\eqref{eq:einstein}:
\begin{equation}
    E = mc^2
    \label{eq:einstein}
\end{equation}

The sample data are listed in Table~\ref{tab:sample}.
\begin{table}[h]
    \centering
    \caption{Sample data}
    \label{tab:sample}
    \begin{tabular}{cc}
        \hline
        alpha & 2 \\
        beta  & 3 \\
        \hline
    \end{tabular}
\end{table}

\end{document}
\end{lstlisting}

\subsection{构建与验证}
\begin{lstlisting}[language=bash]
# 安装依赖（首次需要；poppler-utils 提供 pdftotext 用于验证）
sudo apt update
sudo apt install -y texlive-latex-base texlive-latex-recommended texlive-latex-extra latexmk poppler-utils

# 构建 PDF（latexmk 自动重跑，保证交叉引用正确）
latexmk -pdf -halt-on-error report.tex
ls -l report.pdf

# 验证：PDF 中不应出现未解析引用 "??"（无输出即通过）
pdftotext report.pdf - | grep -n '??'
\end{lstlisting}

\subsection{运行结果}
\begin{itemize}
\item report.pdf 成功生成（1 页，约 62 KB），\lstinline|ls -l report.pdf| 确认存在；
\item latexmk 自动重跑 2 次编译，正文中公式引用 \lstinline|\eqref{eq:einstein}| 显示为 (1)，表格引用 \lstinline|\ref{tab:sample}| 显示为 1；
\item \lstinline|pdftotext report.pdf - | grep -n '??'| 无输出，检查脚本输出 \lstinline|OK: PDF 中没有 ??|、\lstinline|OK: 日志中无 undefined references|。
\end{itemize}

\subsection{要点分析}
\begin{itemize}
\item \textbf{公式编号与引用}：\lstinline|\begin{equation} ... \label{eq:einstein} \end{equation}| 自动编号，正文用 \lstinline|\eqref{eq:einstein}| 引用。
\item \textbf{表格编号与引用}：\lstinline|\caption{Sample data}| 生成标题，\lstinline|\label{tab:sample}| 绑定编号，正文用 \lstinline|\ref{tab:sample}| 引用；表格为严格的 2 行 2 列。
\item \textbf{交叉引用机制}：引用信息写入 .aux 文件，需要编译两遍才能解析；latexmk 会自动重跑直到引用稳定，因此最终 PDF 不会出现 ``??''。
\item \textbf{中文支持}：模板为英文内容，pdflatex 可直接编译；如文档中需要使用中文，可安装 texlive-lang-chinese，用 \lstinline|\usepackage[UTF8]{ctex}| 配合 \lstinline|latexmk -xelatex| 编译（见 scripts/report-zh.tex）。
\end{itemize}

\section{遇到的问题与解决}
\begin{center}
\begin{tabular}{@{}p{5cm}p{5.2cm}p{5cm}@{}}
\toprule
问题 & 原因 & 解决 \\
\midrule
新环境 git commit 报 “Please tell me who you are” & 未配置 user.name / user.email & 提交前执行 git config user.name、git config user.email \\
验证命令 pdftotext: command not found & pdftotext 属于 poppler-utils 软件包，未安装 & apt 安装 poppler-utils \\
交叉引用初次编译显示 “??” & 引用解析依赖 .aux 文件，需编译两遍 & 使用 latexmk，其会自动重跑直至引用稳定 \\
\bottomrule
\end{tabular}
\end{center}

\section{实验总结}
通过本次实验，我掌握了含空格文件名的安全处理方式（\lstinline|find -exec| 传参 + \lstinline|cp --parents| 保留目录结构）与目录/文件权限设置（750/640）；熟练运用 \lstinline!tail | awk | sort | uniq -c | sort | head! 管道完成日志统计，并手算验证了统计结果；理解了 Git 合并冲突的产生条件（双方从同一基线修改同一行）与标准解决流程（编辑、add、commit）；掌握了 LaTeX 中公式、表格的自动编号与交叉引用（label / ref / eqref），以及 \lstinline|latexmk -pdf -halt-on-error| 的命令行构建与验证方法。

\section*{附录：文件清单}
\begin{lstlisting}
scripts/q01.sh         第 1 题完整脚本
scripts/q02.sh         第 2 题完整脚本（含 access.csv 与 analyze.sh 的生成）
scripts/q03.sh         第 3 题完整脚本
scripts/q04.sh         第 4 题完整脚本（含依赖安装、构建、验证）
scripts/report.tex     第 4 题 report.tex（英文版）
scripts/report-zh.tex  第 4 题中文版 report.tex（需 texlive-lang-chinese + xelatex）
实验报告.tex           本报告 LaTeX 源码（ctex + xelatex 编译）
实验报告.pdf           编译生成的报告 PDF（8 页）
\end{lstlisting}

\end{document}
